% =====================================================================
\subsection{ETSI GS QKD 002: the use cases}
\label{sec:qkd002}
% =====================================================================

GS~QKD~002 V1.1.1 (June 2010)~\cite{qkd002} is the first document published
by the ISG QKD. Its goal is not strictly technical: before standardising
\emph{how} to implement QKD, the group wanted to clarify \emph{where} it is
needed, \emph{who} uses it and \emph{which requirements} it must meet to
enter existing infrastructures. The document reasons in three steps: it
places QKD among cryptographic primitives, identifies the layers of the
network stack where it can be inserted and describes six use cases with
their requirements. Following this reasoning is useful because it shows
where the needs come from that, twelve years later, led to GS~QKD~015
and~018.

\subsubsection{QKD as a cryptographic primitive}

The document starts from a typical QKD system: two units, usually 19"
rack-mounted, connected by a quantum channel over fibre or free space up to
about $100$~km long, which output the same random bit sequence at both
ends. Clause 4 immediately makes clear that this solves \emph{a single}
cryptographic problem.

A secure system needs three primitives: \emph{encryption}, \emph{key
distribution} and \emph{message authentication}. QKD covers only the
second; encryption requires a one-time pad or symmetric ciphers such as
AES, and authentication requires MACs (for instance based on Wegman–Carter
universal hashing) or digital signatures. It follows that QKD is a
\emph{complement} to classical cryptography, not a replacement, and that
its overall security also depends on the other components of the system.
The document also stresses that the public channel of QKD must be
authenticated with a pre-distributed secret.

\subsubsection{Where QKD fits in the network stack}

Clause 6 describes at which layer of the stack a QKD key can be consumed.
The choice of layer determines what must be modified and how transparent
the integration is.

\begin{description}[style=nextline,leftmargin=1.2em]
  \item[Data link layer (L2)] The QKD key feeds the PPP key exchange, the
        IEEE~802.1AE (MACsec) keys or a \emph{QKD link encryptor}, i.e.\ a
        line encryptor that directly consumes the keys produced by the link
        and creates an encrypted tunnel between two sites. It is the most
        transparent option for applications.
  \item[Network layer (L3)] A modified IKE uses the QKD key as the shared
        secret for IPsec.
  \item[Transport layer (L4)] The QKD key replaces the TLS/SSL session keys
        or is used instead of HMACs for authentication.
  \item[Application layer] Applications use the key to authenticate users
        or to encrypt the content exchanged between their instances.
\end{description}

The conceptual point is that, whatever the layer, there is one component
that \emph{consumes} keys and one that \emph{supplies} them: this is the
separation that later becomes SAE and KME, with a standard interface
(GS~QKD~004 and~014) between them.

\subsubsection{The six use cases}

Clause 7 describes six use cases with a fixed structure: goal,
description, actors, benefits, operational and quality-of-service
considerations. We summarise them along a logical thread: from a single
link to a network, and from a private network to an operator's network.

\paragraph{UC1 – Offsite backup / business continuity.}
An enterprise replicates data from a primary site to a backup site (daily
backups, snapshots, SAN mirroring, database replication). A single
point-to-point QKD link feeds a line encryptor over Ethernet or Fibre
Channel; the fibre may be dedicated (\emph{dark fibre}) or shared via WDM
or TDM. The requirements that emerge are: key renewal as frequent as the
generation rate allows, handling of recalibration phases during continuous
transfers, an initial secret at installation, and high encryption speed.
The document notes that overall security is not unconditional, since it
also depends on the encryptor.

\paragraph{UC2 – Enterprise metropolitan area network.}
Several branches of an enterprise or a public body are connected to data
centres, through which branch-to-branch traffic also flows. Each link is
protected by a QKD link encryptor that encrypts and authenticates at
layer~2, transparently for applications. Two new needs appear: different
links carry different traffic volumes and timings, hence \emph{different
QoS requirements}, and transfers may be scheduled, periodic or triggered by
thresholds. The \emph{QKD Secured Key Server} variant has a central server
that distributes keys and certificates to clients across the metropolitan
area over channels encrypted with one-time pad and authenticated with
Wegman–Carter: an early example of a key service offered by a network.

\paragraph{UC3 – Critical infrastructure control and data acquisition.}
This case protects the SCADA traffic of a critical infrastructure, with a
railway network as the example: here the \emph{authenticity} of commands
(switches, signals) matters more than their confidentiality. The nodes of
the private WAN sit in stations with a controlled perimeter, while the
fibre lies outside it. The document proposes two architectures: single QKD
links between nodes, with key buffering and forwarding handled in the
nodes, or a \emph{QKD network with integrated key management} and a
dedicated interface towards applications. It asks for the highest
availability, with precautions for recalibration, failures and denial of
service, and for the ability to initialise new redundant links with secrets
obtained along an alternative path. It is the first case in which key
relay, centralised management and path redundancy appear together.

\paragraph{UC4 – Backbone protection.}
A fibre in an operator's metropolitan backbone, typically a WDM ring with
OADMs or ROADMs, carries a quantum channel alongside the data channels. The
keys protect control and management traffic or are offered to customers.
There are two actors: the infrastructure provider and the service provider
that uses it; backbone nodes are trusted. The requirements are the closest
to those of an operator: the QKD system must not reduce the throughput of
the WDM system by more than $10\%$, must be compatible with common fibres,
splices and connectors, its management must \emph{integrate into the
management of the WDM system}, and installation and maintenance must be
possible for staff without quantum physics expertise.

\paragraph{UC5 – High-security access network.}
In a passive optical network (PON) for Fiber-to-the-Home, downstream traffic
is broadcast to all ONUs, so it must be encrypted to prevent one user from
reading another's data. One OLT serves $32$ to $128$ ONUs and the quantum
channel follows the same path as the classical one. The configuration is
highly asymmetric: a single, more expensive detector at the OLT serves many
sources at the ONUs. Short distances allow high rates, which are needed
because the generation capacity is shared among all users. The \emph{QKD
Authenticated Sensor Network} variant applies the same scheme to a central
station with many sensors and actuators.

\paragraph{UC6 – Long-haul service.}
Two distant ground stations $A$ and $B$, for example on opposite shores of
an ocean, want to share a key without trusting intermediate ground nodes. A
moving node $C$ (satellite, aircraft or high-altitude platform) establishes
a secret $a$ with $A$ and then a secret $b$ with $B$, and sends $a\oplus b$
to $B$ over a classical channel; $B$ recovers $a$. It is a key relay in
which the trusted node is $C$. Low-orbit satellites fly at about
$300$–$800$~km and pass over the stations about once a day; geostationary
orbit is not practical. Weather and atmosphere can reduce or cancel the
rate, so redundant channels are needed.

Table~\ref{tab:qkd002} summarises the six cases.

\begin{table}[h]
\centering\small
\caption{The GS~QKD~002 use cases and the elements that anticipate later
specifications.}
\label{tab:qkd002}
\begin{tabularx}{\textwidth}{@{}l L L@{}}
\toprule
\textbf{Case} & \textbf{Configuration} & \textbf{Need it anticipates} \\
\midrule
UC1 Backup & point-to-point link + line encryptor & key–encryptor interface \\
UC2 Enterprise metro & several links, key server & per-application QoS \\
UC3 Critical infrastructure & network with key relay and integrated management & KMS, redundancy, centralised control \\
UC4 Backbone & quantum channel in WDM & coexistence with the optical network, integrated management \\
UC5 PON access & point-to-multipoint & sharing of key capacity \\
UC6 Long-haul & relay through a moving node & trusted node, redundancy \\
\bottomrule
\end{tabularx}
\end{table}

\subsubsection{What it leaves behind}

GS~QKD~002 is a first-edition document: the clause on general requirements
and the functional characteristics of the individual cases are marked as
not yet available. Its value lies in having set the questions. Read in
sequence, the use cases show that useful QKD is not a single link but a
\emph{network}, which requires: a key management system that stores and
relays keys (UC3, UC6); a way to express different service requirements for
different applications (UC2); and coexistence with the operator's optical
network, with management integrated into the existing one and usable by
non-specialist staff (UC4). The first two needs are answered by GS~QKD~015,
the third, in the form of coordination between the QKD network and the
transport network, by GS~QKD~018.
